% Generated draft: requires mathematical and source review.
\documentclass[11pt]{article}
\usepackage[margin=25mm]{geometry}
\usepackage{fontspec}
\setmainfont{Latin Modern Roman}
\usepackage{amsmath,amssymb,mathtools,unicode-math}
\setmathfont{Latin Modern Math}
\usepackage{enumitem,needspace,xurl,hyperref}
\setlength{\parindent}{0pt}
\setlength{\parskip}{5pt}
\setlength{\emergencystretch}{3em}
\newcommand{\N}{\mathbb{N}}
\newcommand{\Z}{\mathbb{Z}}
\newcommand{\Q}{\mathbb{Q}}
\newcommand{\R}{\mathbb{R}}
\newcommand{\C}{\mathbb{C}}
\newcommand{\F}{\mathbb{F}}
\newcommand{\A}{\mathbb{A}}
\newcommand{\PP}{\mathbb{P}}
\newcommand{\E}{\mathbb{E}}
\newcommand{\eps}{\varepsilon}
\newcommand{\dd}{\,\mathrm{d}}
\begin{document}
\section*{UnsolvedMath ID: 30001312}
\textbf{Source problem code:} catalogue entry
\subsection{Definitions and mathematical statement}
Let \(\Sigma\) denote the surface intended in the source; its topology, compactness, genus, and boundary conditions are not specified in the supplied record. For a Lie group \(G\), write
\[
\mathcal X_G(\Sigma)=\operatorname{Hom}(\pi_1\Sigma,G)/G,
\]
where \(\operatorname{Hom}(\pi_1\Sigma,G)\) is the set of group homomorphisms and \(G\) acts by conjugation. This notation specifies conjugacy classes; the supplied record does not specify a topology or an algebraic quotient convention. Here \(\operatorname{SU}(2,1)\) is the determinant-one group preserving a nondegenerate Hermitian form of signature \((2,1)\), and \(\operatorname{PSL}(2,\C)=\operatorname{SL}(2,\C)/\{\pm I\}\).

Let \(\mathbb{CH}^{2}\) be complex hyperbolic two-space with its Kähler metric and Kähler form \(\omega\). A real surface immersion \(f\) into \(\mathbb{CH}^{2}\) is minimal Lagrangian if its mean-curvature vector vanishes and \(f^*\omega=0\). In an equivariant formulation, a representation \(\rho:\pi_1\Sigma\to\operatorname{SU}(2,1)\) and an immersion \(f:\widetilde\Sigma\to\mathbb{CH}^{2}\) satisfy
\[
f(\gamma x)=\rho(\gamma)\cdot f(x)
\qquad(\gamma\in\pi_1\Sigma,\ x\in\widetilde\Sigma).
\]
This describes the usual equivariance condition, but the record does not identify the full class of solutions or its equivalence relation.

\textbf{Source conjecture, with unresolved scope.} Do the minimal Lagrangian surface solutions under consideration parametrize a portion of
\[
\mathcal X_{\operatorname{SU}(2,1)}(\Sigma)
\]
analogous to the quasi-Fuchsian component of
\[
\mathcal X_{\operatorname{PSL}(2,\C)}(\Sigma)?
\]
For context only, when \(\Sigma\) is a closed oriented surface of genus at least two, the classical quasi-Fuchsian locus consists of conjugacy classes of discrete faithful representations into \(\operatorname{PSL}(2,\C)\) whose limit set in the Riemann sphere is a Jordan curve. That hypothesis on \(\Sigma\) is not supplied here and is not imposed on the conjecture.

The source does not define the proposed solution parameter space, the target portion, the required meaning of ``parametrize,'' or the properties constituting the quasi-Fuchsian analogy. Consequently, no definite bijection, homeomorphism, or global parametrization theorem can be extracted from this record alone.
\subsection{Short English statement}
The conjecture asks whether the intended minimal Lagrangian surfaces in complex hyperbolic two-space describe a region of surface-group representations into \(\operatorname{SU}(2,1)\), analogous to the classical quasi-Fuchsian region. The surface hypotheses, solution space, target region, and precise meaning of parametrization remain unspecified.
\subsection{Sources}
\begin{itemize}
\item Local UnsolvedMath registry: \texttt{data/unsolvedmath/problems.json}, permanent ID \(30001312\), exact source problem code \texttt{catalogue entry}, title ``Complex-Hyperbolic Quasi-Fuchsian Representation Components.'' The supplied record belongs to the Oberwolfach Reports Open Problems collection and gives proposed year \(2009\); no individual proposer is supplied.
\item Supplied report citation: \emph{Differentialgeometrie im Großen (2010)}, Oberwolfach Reports, pp.~1943--1972, DOI \texttt{10.4171/owr/2009/35}. Source URL and literature-source link labelled ``Differential geometry report'': \url{https://doi.org/10.4171/owr/2009/35}. The supplied background describes Oberwolfach Reports as documenting mathematical workshops and their open-problem discussions.
\item Supplied literature-source link labelled ``McIntosh--Loftin theorem'': \url{https://doi.org/10.1007/s10711-012-9717-1}. The supplied record attributes to McIntosh--Loftin the construction of a neighborhood of \(\R\)-Fuchsian representations and a proof that this neighborhood is complex-hyperbolic quasi-Fuchsian; no fuller bibliographic citation or theorem hypotheses are supplied.
\item Literature assessment dated exactly \texttt{2026-08-21}: the supplied curation labels the problem ``partially solved'' and states that the conjectured global parametrisation of the representation space remains open. This is supplied dated curation derived from the OWR material, not an independently verified literature assessment. Its time-sensitive status and cited claims require independent verification.
\item The supplied registry labels the statement \texttt{corrected\_verified} and reports that its mathematical content was checked against the official Oberwolfach report, with notation, context, and question form normalized without changing the source-level claim. This is a supplied verification assertion; no independent source consultation is claimed here.
\end{itemize}
\end{document}
